\section{Related Work}
\label{submission:related}

\textbf{Model Merging and Parameter Averaging.}
The simplest and most direct form of combining models is parameter averaging, popularized by Model Soups \cite{wortsman2022robust}, which averages the weights of models fine-tuned with different hyperparameters or random seeds.
To merge models trained on different initializations or architectures, researchers have explored aligning the permutation symmetries of neural network channels.
\emph{Git Re-Basin} \cite{ainsworth2022git} formulates this as a linear assignment problem to align neurons before averaging.
Similarly, \emph{REPAIR} \cite{jordan2022repair} rescales activation statistics to mitigate the loss of variance (the ``barrier'' phenomenon) during interpolation, while \emph{ZipIt!} \cite{stoica2023zipit} merges features by matching and nesting overlapping representations.
While these methods are conceptually interesting, they either require expensive permutation search or are limited to combining models with identical or structurally aligned backbones.

\textbf{Task Vectors and Task Arithmetic.}
To enable modular model editing, Ilharco et al.~\cite{ilharco2022editing} introduced the concept of \emph{task vectors}, which are formed by subtracting the pre-trained base model weights from the fine-tuned expert weights.
This formulation allows task-specific knowledge to be added, subtracted, or scaled via standard vector arithmetic in weight space.
Task Arithmetic has been successfully applied to steering language model generation, modifying image-text retrieval styles, and consolidating specialized experts.
However, simple linear addition of task vectors assumes that parameters in different models are independent and orthogonal.
In practice, this assumption breaks down due to high parameter overlap, leading to severe destructive interference when multiple models are merged or when large scaling coefficients are used.

\textbf{Conflict Resolution in Weight Space.}
To resolve the destructive interference inherent in Task Arithmetic, several weight-space conflict resolution methods have been proposed.
\emph{TIES-Merging} \cite{yadav2023ties} addresses sign conflicts by identifying that redundant parameter updates can be pruned and that a sign consensus can be computed via voting.
However, TIES-Merging requires setting a trimming threshold ($k\%$), executing sign voting, zeroing out non-conforming parameters, and then scaling the merged vector back up to match the original energy.
This multi-stage pipeline is complex and heavily dependent on finding optimal hyperparameter values.
To reduce this complexity, \emph{DARE} \cite{yu2024dare} uses randomized dropout to sparsify task vectors, but still requires scaling parameters to offset the information loss and introduces stochastic noise.
Other lines of work, such as \emph{AdaMerging} \cite{yang2023adamerging} and \emph{SyMerge} \cite{yang2024symerge}, optimize merging coefficients (either task-wise or layer-wise) on unlabeled test data using gradient descent.
While effective, these adaptive methods require test-time optimization, violating the training-free promise of model merging.
Mathematical alignment methods like \emph{OrthoMerge} and \emph{SAIM} apply Riemannian manifold mappings or Singular Value Decomposition (SVD) to preserve weight geometry, but their cubic computational complexity ($O(D^3)$ where $D$ is the hidden dimension) renders them highly impractical for modern large-scale neural networks.

In contrast to these increasingly convoluted pipelines, our proposed WTA-Sign method aligns with the principles of \emph{minimalist deep learning}.
We show that by using the magnitude of parameter updates as a proxy for task confidence, we can eliminate the entire trimming, voting, and scaling pipeline of TIES-Merging, as well as the optimization loops of adaptive merging, while achieving superior multi-task performance with zero hyperparameters and exceptional computational efficiency.
